
\sourcenote{Genre: admissible rule (tier 2) --- a proof-tree-rewriting
recipe. The closure discipline this move belongs to --- why strict
literals never reach a row, and how the three moves divide the
positions --- is set out in full in the
\rulename{ClosureDiscipline-EXACT} section; this section states only
what is native to the outward move: the two verified faces of its
warrant, and the direction analysis that confines it to negative
positions.}

\begin{brecipe}[the admissible template, general form]\label{rec:co-template}
An admissible rule $R$ is warranted by a rewrite replacing any
$R$-node with a derivation fragment over core (and
previously-justified) rules with the same conclusion --- or, where the
conclusion has no catalogue derivation, discharging against a named
verified theorem. Leaf-first application eliminates every $R$-node in
finitely many steps; soundness is conserved with no new kernel
semantic content.
\end{brecipe}

\begin{brecipe}[the \rulename{ClosureDiscipline-OUTWARD} instance --- two faces]
\label{rec:co-instance}
\emph{Atom face.} At a negative position, the complement of a closed
atom is consumed as its orientation swap:
\[
\mathrm{cmpl}\{p \le c\} \;\rightsquigarrow\; \{c \le p\}
\qquad\qquad
\mathrm{cmpl}\{c \le p\} \;\rightsquigarrow\; \{p \le c\}
\]
warranted by the containment
$\mathrm{cmpl}\{p \le c\} \subseteq \{c \le p\}$
(Proposition~\ref{prop:co-atom}): the consumed atom
\emph{over-approximates} the true complement, losing only the cut
$\{p = c\}$.

\emph{Region face.} A negative position consuming the complement of a
\emph{certified region} $R$ (a child certificate\textquotesingle s
conclusion, not a syntactic atom) is warranted by the antitone kernel
(Proposition~\ref{prop:co-region}): certified under-approximation of
the denotation becomes over-approximation of the
denotation\textquotesingle s complement. Both faces discharge against
their verified containment; the orientation swap itself is syntactic
and recomputed by the checker.
\end{brecipe}

\begin{bprop}[atom face]\label{prop:co-atom}
$\mathrm{cmpl}\{x \mid p(x) \le c\} \subseteq \{x \mid c \le p(x)\}$,
and symmetrically
$\mathrm{cmpl}\{x \mid c \le p(x)\} \subseteq \{x \mid p(x) \le c\}$.
\end{bprop}
\begin{proof}
If $\neg(p(x) \le c)$ then $c < p(x)$, hence $c \le p(x)$ --- the
strict fact weakened to its non-strict form
(\texttt{closureOutward\_sound}; the mirror is
\texttt{closureOutwardGe\_sound}). The containment is strict exactly
on the cut $\{p = c\}$, which the swap adjoins.
\end{proof}

\begin{bprop}[region face --- the antitone kernel]\label{prop:co-region}
If a certified region under-approximates a denotation,
$R \subseteq [\![\varphi]\!]$, then
$\mathrm{cmpl}\,[\![\varphi]\!] \subseteq \mathrm{cmpl}\,R$: the
complement of the certificate over-approximates the complement of the
denotation.
\end{bprop}
\begin{proof}
Contrapositive, one line: $x \notin \mathrm{cmpl}\,R$ means
$x \in R \subseteq [\![\varphi]\!]$, so
$x \notin \mathrm{cmpl}\,[\![\varphi]\!]$
(\texttt{compl\_over\_approx}). This is the direction transport that
lets negative positions consume child certificates at all: the
judgement kind fixes one orientation, and complementation flips it in
precisely the way the consuming side wants.
\end{proof}

\begin{bremark}[what is lost, and why the direction is right]
Only the cut $\{p = c\}$ --- and in the sound direction. At a
negative position the atom names a set to be excluded or avoided;
certifying against a superset of it is a fortiori certifying against
it. The move never claims the cut is negligible --- for a nondegenerate
observable ($p - c \not\equiv 0$ on a full-dimensional cell) it is a
measure-zero level set, but soundness nowhere consumes that: the
containment alone carries the warrant, so even a degenerate observable
whose cut is fat (at the extreme $p \equiv c$, the whole cell) only
makes the certificate more conservative, never wrong.
\end{bremark}

\begin{bremark}[one-directional by design]
At a \emph{positive} position the same swap would overclaim. Take $p$
constant at $c$: the swap $\{c \le p\}$ is the whole space, while the
true complement of $\{p \le c\}$ is empty --- the maximal over-claim.
Nothing in the move\textquotesingle s own premises excludes this; the
closure discipline\textquotesingle s routing table
(the preamble, \rulename{ClosureDiscipline-EXACT} section) is the
guard, sending positive-position strict literals to the inward move
instead. A porter implementing the discipline ports the routing with
it, or re-learns this remark the hard way.
\end{bremark}

\begin{bremark}[what is deliberately not warranted here]
The verified faces are the atom step and the region step. The
composite induction --- pushing complements through conjunctions,
disjunctions, and nested operators until they reach atoms --- is
stated once, for all consumers, with the normal-form rewriting rules,
and this rule\textquotesingle s dashboard check is marked conditional
for exactly that reason: a complement traversing a full specification
rides the normal-form sections for the traversal and this rule for
the atoms it arrives at.
\end{bremark}

\begin{anchors}
Proposition \ref{prop:co-atom} & \texttt{closureOutward\_sound}, \texttt{closureOutwardGe\_sound}\\
Proposition \ref{prop:co-region} & \texttt{compl\_over\_approx}\\
the judgement reading & \texttt{RegionSat} (quoted in full above)\\
the routing guard & the closure-discipline preamble (\rulename{ClosureDiscipline-EXACT} section)\\
the composite traversal & the normal-form rewriting sections (deliberately not this rule)\\
\end{anchors}
